\section{Reconstruction and the actual HYM flow}\label{sec:actual-flow}

On the actual rank-four bundle we use physical time $s$ and the equation
\begin{equation}\label{eq:hym-flow}
 h^{-1}\partial_sh=-2i\Lambda_{\omega_t}F_h.
\end{equation}
The Einstein constant is zero. All base and bundle parameters are fixed
in this equation. Relative to $h_0$, all metrics under consideration
have determinant one. For two rank-$r$ metrics put
\begin{equation}\label{eq:trace-distance}
 \sigma_r(h,k)=\tr(h^{-1}k+k^{-1}h)-2r.
\end{equation}
This expression is intrinsic, nonnegative, and at least $d_r(h,k)^2$.
On a fixed positive metric range it is uniformly equivalent to that
squared distance. For an approximate metric $k$, write
\[
 \mathcal R(k)=k^{-1}\partial_sk+2i\Lambda_{\omega_t}F_k.
\]

\begin{theorem}\label{thm:actual-tracking}
For each effective trajectory of \eqref{eq:finite-flow} with initial tuple
in a fixed bounded set, and all sufficiently small $t$, there is a smooth
determinant-one comparison metric $\widehat K(s)$ on $\cB_{t,\theta}$,
defined for every $s\ge0$, with uniform positive bounds against $h_0$,
such that
\begin{equation}\label{eq:global-residual}
 \sup_{s\ge0}\int_{X_t}|\mathcal R(\widehat K(s))|_{\op}\,dV_t
        \le C\frac{(1+\log L)^2}{L^2}.
\end{equation}
Let $h(s)$ solve \eqref{eq:hym-flow} from any smooth determinant-one
metric with $ch_0\le h(0)\le Ch_0$, for fixed positive bounds.
Set $\eta=\int_{X_t}\sigma_4(h(0),\widehat K(0))\,dV_t$ and
$e_L=(1+\log L)^2/L$. Then, for every $T\ge0$,
\begin{equation}\label{eq:bounded-tracking}
 \sup_{0\le s\le TL}\int_{X_t}d_4(h(s),\widehat K(s))^2\,dV_t
          \le\eta+CTe_L.
\end{equation}
The constants are independent of $\theta$, $t$, $s$ and $T$ in their
stated domains. The exact initial metric need not be a tensor metric.
\end{theorem}

\subsection{Equivariant reconstruction on the elliptic family}
The reconstruction uses a component correction cancelling the leading
velocity, flat geodesic interpolation on the necks, and an estimate
for their mismatch on the collars. We then tensor with the fixed
K3 metric and descend, keeping the full contracted Chern equation.
Put $\epsilon=L^{-1}$ and choose mean-zero component Green functions
\begin{equation}\label{eq:greens}
 \Delta_jg_{je}=-2\pi\delta_{p_{je}}+\frac{2\pi}{A_j},\qquad
 g_{je}=\log|u|+\gamma_{je}(u),
\end{equation}
where $\Delta_j$ is the nonnegative scalar Laplacian on $D_j$ and
$u$ is its node coordinate. Uniqueness makes these functions equivariant
under the isometry $\sigma$. The two actual incoming transports to
component $j$ are $U_{j,j-1}=R$ and $U_{j,j+1}=R^*$.
Define
\begin{equation}\label{eq:component-correction}
 B_{je}=\log(H_j^{-1}U_{jk}H_kU_{jk}^*),\qquad
 \xi_j=-\sum_{e\ni j}B_{je}g_{je},\qquad K_j=H_j e^{\epsilon\xi_j}.
\end{equation}
The logarithms are $H_j$-self-adjoint and trace-free. On the punctured
component, the effective equation gives the exact cancellation
\begin{equation}\label{eq:component-cancellation}
 \epsilon\Delta_j\xi_j=-H_j^{-1}\dot H_j.
\end{equation}

Cut at radius $\rho=L^{-1}$. At incidence $je$ write
\[
 \xi_j=-B_{je}\log|u|+c_{je}+q_{je}(u),\qquad q_{je}(0)=0.
\]
The corrected endpoint is $P_{je}=H_je^{\epsilon c_{je}}$.
On the node $j,k$, in the $j$ incidence frame, use
\begin{equation}\label{eq:node-geodesic}
 N_e(u)=\Geo(P_{je},U_{jk}P_{ke}U_{jk}^*;-\epsilon\log|u|).
\end{equation}
Reversing the geodesic and using $uv=t$ gives the same metric in the
other incidence frame. For each fixed time this node metric has zero
Chern curvature: if $P,Q$ are fixed endpoints, it is
$P\exp(-\epsilon\log|u|\log(P^{-1}Q))$, whose logarithmic derivative
has zero $\bar\partial$ away from the puncture.

Interpolate relative logarithms between $K_j$ and $N_e$ on
$\rho<|u|<2\rho$. Choose paired cutoffs, component transfers and node
charts under the involution. On fixed components the required smooth
transfer maps can be obtained from invariant lifts of base vector fields.
No holomorphic product trivialization is assumed. The operations
in \eqref{eq:component-correction}--\eqref{eq:node-geodesic} commute
with unitary transport, geodesic reversal and the actual lift
\eqref{eq:projective-lift}. Thus the resulting curve metric $K_E(s)$
is invariant in component \emph{and} node frames.

We give the residual estimate for this full Chern construction.
In a flat holomorphic frame, with
$\omega_{E_t}=(i/2)a\,dz\wedge d\bar z$, the actual equation is
\begin{equation}\label{eq:curve-chern-flow}
 \partial_sH=4a^{-1}(H_{z\bar z}-H_{\bar z}H^{-1}H_z).
\end{equation}
In real coordinates it contains
$i(H_xH^{-1}H_y-H_yH^{-1}H_x)$ as well as the two symmetric quadratic
terms. This commutator is retained in all curvature estimates below.

Effective-flow confinement and speed give, uniformly for all $s\ge0$,
\begin{equation}\label{eq:green-bounds}
 \begin{gathered}
 |B_{je}|\le C,\qquad |\dot B_{je}|\le C\epsilon,\qquad
 \|\xi\|_\infty\le C(1+|\log\rho|),\\
 \int\bigl(|\xi|+\epsilon^{-1}|\dot\xi|+|\nabla\xi|\bigr)\le C,
 \qquad \int|\nabla\xi|^2\le C(1+|\log\rho|).
 \end{gathered}
\end{equation}
These integrals are over the truncated components with their fixed metrics.
The identity
\[
 e^{-h}\partial e^h=\int_0^1e^{-uh}(\partial h)e^{uh}\,du
\]
and its differentiated form show that the component residual, after
\eqref{eq:component-cancellation}, is bounded by the integrals of
\[
 C\epsilon^2\bigl(|\xi|+|\xi|\,|\Delta_j\xi|
                                 +|\nabla\xi|^2\bigr).
\]
The second derivative is contracted before taking its norm: in complex
dimension one the $\bar\partial\partial\xi$ term in each conjugation
is a scalar multiple of $\Delta_j\xi$. Differentiating the conjugating
factors gives the displayed quadratic-gradient terms. Since
$\Delta_j\xi$ on the punctured component is bounded, the residual is
$O(\epsilon^2(1+|\log\rho|))$ after integration.
The time derivative of the correction has the same bound by
\eqref{eq:green-bounds} and $H_j^{-1}\dot H_j=O(\epsilon)$.

Here is the collar estimate which allows matching without losing this
order. For small endpoint logarithms $C_j,C_k$, let $D(C_j,C_k;u)$
be the relative logarithm at $H_j$ of
$\Geo(H_je^{C_j},U_{jk}H_ke^{C_k}U_{jk}^*;u)$, minus
$C_j+uB_{je}$. On the confined matrix range it is smooth, and
\[
 D(0,0;u)=0,\qquad D(C_j,C_k;0)=0.
\]
Bounded mixed derivatives imply
$|D|\le C(|C_j|+|C_k|)|u|$ and
$|\partial_uD|+|\partial_u^2D|\le C(|C_j|+|C_k|)$.
Substitute $C_j=\epsilon c_{je}$, $C_k=\epsilon c_{ke}$ and
$u=-\epsilon\log|z|$. Including the regular part $q_{je}$ gives
the mismatch bound, through two spatial derivatives,
\begin{equation}\label{eq:collar-mismatch}
 C\bigl(\epsilon\rho^{1-k}
             +\epsilon^2(1+|\log\rho|)\rho^{-k}\bigr),
 \qquad k=0,1,2.
\end{equation}
In the second derivative the additional term of order
$\epsilon^3\rho^{-2}$ is absorbed by the displayed bound.

The curvature difference from the exactly flat node profile contains
the second derivative of this mismatch, its first derivative times the
node logarithmic derivative $O(\epsilon/\rho)$, and its zeroth
derivative times the corresponding second-derivative and quadratic
terms. Cutoff derivatives have the same orders in $\rho$.
Since the collar area is $O(\rho^2)$, its integrated residual is
\[
 C\bigl(\epsilon\rho+
                 \epsilon^2(1+|\log\rho|)^2\bigr).
\]
The middle-node time residual is $O(\epsilon\rho^2)$; its spatial
Chern curvature is zero. Transfer to the actual smooth fibre adds
$C|t|\rho^{-N}$ for a fixed finite $N$, by the smooth plumbing
metric coefficients and the two derivatives in the contracted equation.
Putting $\rho=L^{-1}$ proves
\begin{equation}\label{eq:curve-residual}
 \sup_{s\ge0}\int_{E_t}
 |K_E^{-1}\dot K_E+2i\Lambda_{E_t}F_{K_E}|_{\op}\,dA
       \le C\frac{(1+\log L)^2}{L^2}.
\end{equation}
Every time-dependent input was a confined tuple or its
$O(L^{-1})$ velocity. Thus this is a residual bound on the whole
physical half-line, including reducible parameters.

\subsection{Exact tensor curvature and descent}
On the cover set $\widetilde K=h_S\otimes K_E$. Product Chern
connections give exact identities
\begin{equation}\label{eq:product-curvature}
 \begin{split}
 D_{\widetilde K}&=D_S\otimes I+I\otimes D_{K_E},\\
 F_{\widetilde K}&=F_{T^{1,0}S}\otimes I+I\otimes F_{K_E},\\
 \mathcal R(\widetilde K)&=I_{T^{1,0}S}\otimes
       (K_E^{-1}\dot K_E+2i\Lambda_{E_t}F_{K_E}).
 \end{split}
\end{equation}
The last equality uses the zero mean curvature of $h_S$. Equivariance
proved above makes $\widetilde K$ descend to $\widehat K$.
Its determinant is one in the fixed trivialization.
The operator norm of the tensor residual is unchanged, whereas
integration divides by the degree two of the cover. Consequently
\eqref{eq:curve-residual} implies \eqref{eq:global-residual}.

For two tensor metrics there is also the exact distance identity
\begin{equation}\label{eq:tensor-distance}
 d_4(h_S\otimes H,h_S\otimes K)^2=2d_2(H,K)^2.
\end{equation}
Together with \eqref{eq:k3-volume}, division by the cover degree
makes their global quotient squared distance equal to the curve
integral. Neither tensor identity is assumed for an arbitrary exact
solution $h(s)$.

\subsection{Intrinsic comparison for arbitrary rank-four metrics}
On each fixed compact smooth fibre, the classical HYM metric flow
exists for all time; see \cite[Section 4.1, after (4.8)]{hkkp}.
Order comparison \cite[Proposition 4.3]{hkkp}
with $ch_0$ and $Ch_0$ preserves
the bounds, and the scalar determinant equation preserves the
normalization. These fixed-fibre statements impose no uniform
Sobolev-constant hypothesis on the degenerating family.

For an exact flow $h$ and a comparison metric $k$ in a common fixed
positive range, the Chern trace calculation gives
\begin{equation}\label{eq:forced-trace}
 (\partial_s+\Delta_{\omega_t})\sigma_4(h,k)
                \le C|\mathcal R(k)|_{\op}.
\end{equation}
For two exact velocities, the intrinsic trace calculation of
\cite[Proposition 2.4]{wang}, with the coupling section set to zero,
makes the left side a sum of nonpositive
quadratic derivative terms in local holomorphic frames. Replacing the
velocity of $k$ by its exact-flow velocity plus $k\mathcal R(k)$
adds the traces of $h^{-1}k\mathcal R(k)$ and its inverse-metric
counterpart. Their absolute values are bounded by the right side
under the fixed relative bounds. The same calculation is summed over
all complex coordinate directions. It is intrinsic and therefore
applies across the nonflat tangent-bundle frames, without reducing
the exact solution to a tensor ansatz.

\begin{proof}[Proof of Theorem~\ref{thm:actual-tracking}]
The equivariant construction and \eqref{eq:product-curvature} prove
the comparison metric and its residual estimate. Integrate
\eqref{eq:forced-trace} on the closed smooth fibre, so that the
Laplacian integrates to zero. Its accumulated right side up to time
$TL$ is at most $CT(1+\log L)^2/L$. Since $d_4^2\le\sigma_4$,
this proves \eqref{eq:bounded-tracking}. No derivative bound on
the exact initial metric enters this argument.
\end{proof}

\begin{theorem}\label{thm:all-time}
Under the hypotheses of Theorem~\ref{thm:actual-tracking}, assume also
$\mu_0(\theta)\ge\kappa>0$. For sufficiently small $t$,
\begin{equation}\label{eq:all-time}
 \sup_{s\ge0}\int_{X_t}d_4(h(s),\widehat K(s))^2\,dV_t
       \le C\eta+C_\kappa e_L(1+|\log e_L|).
\end{equation}
\end{theorem}

\begin{proof}
For the actual flow write $h=h_0H=h_0e^{\Psi}$. Pointwise $\tr \Psi=0$.
The exact relative Chern equation is
\[
 H_s=-D_\theta^*D_\theta H
       -4\sum_\alpha(D_\alpha H)^*H^{-1}(D_\alpha H).
\]
The contracted background-curvature commutator is zero. All quadratic
terms are retained, including those from noncommuting derivatives.
Put $F(s)=\int_{X_t}\tr(e^{\Psi}-I)\,dV_t$. On the fixed positive range,
Taylor's formula and trace freeness give
\[
 \frac c2\|\Psi\|_{L^2}^2\le F\le\frac C2\|\Psi\|_{L^2}^2,
 \qquad -F'\ge\frac{c^2}{C}\int_{X_t}|D_\theta \Psi|^2\,dV_t.
\]
For the derivative inequality first bound $H^{-1}\ge C^{-1}I$
inside the positive quadratic term. The unweighted trace then satisfies
$4\sum_\alpha\tr((D_\alpha H)^*D_\alpha H)=|D_\theta H|^2$.
The divided differences of the exponential on the bounded spectrum
are at least $c$, giving $|D_\theta H|\ge c|D_\theta \Psi|$.
No commutator is discarded inside a matrix-weighted trace.

There is no trace-free parallel kernel on the stated parameter class.
Theorem~\ref{thm:five-spectrum} therefore gives
\[
 F(s)\le F(0)\exp\left(
 -\frac{2c^2}{C^2}(1-\varepsilon(t))\frac{2\pi\mu_0(\theta)}{L}s
 \right).
\]
Lemma~\ref{lem:finite-decay} gives a corresponding exponential tail
for the effective trajectory. The reconstruction is uniformly Lipschitz
in global $L^2$ on the confined tuple range: the Green coefficients
have bounded $L^2$ norms multiplied by $L^{-1}$, and the node and
collar maps have bounded tuple derivatives. It sends the identity tuple
exactly to $h_0$. Equation~\eqref{eq:tensor-distance} therefore
transfers effective-flow decay to the comparison metric.

For $s\ge TL$, the exact and reconstructed squared distances to
$h_0$ are at most $C e^{-c_\kappa T}$. For $s\le TL$, use
\eqref{eq:bounded-tracking}. Taking $T$ proportional to
$1+|\log e_L|$ and applying the squared triangle inequality proves
\eqref{eq:all-time}.
\end{proof}

The all-time conclusion is uniform on a nonempty stable parameter class.
At reducibility there are nonidentity parallel metrics, and as
$\mu_0\to0$ the above tail constants degenerate. The bounded-time
comparison remains uniform there; extending it to the weakest inverse
rate requires a further argument.
